\scanpage{46}
\mspage{\handwritten{14}}
\uline{11.6.}- Soient $G$ un $S$-groupe plat quasi-compact et quasi-séparé sur $S$\handwritten{,} tel que $\handwritten{\mathscr{A}}(G)$ soit un $\handwritten{\mathcal{O}}_S$-Module plat $\handwritten{\mathscr{A}}$, et $\handwritten{\mathcal{E}}$ un $\handwritten{\mathcal{O}}_S$-Module quasi-cohérent et plat. \deleted{Nous dirons qu'}\inserted{\handwritten{Notons qu'}}un morphisme de $S$-foncteurs de $G$ dans $\inserted{\handwritten{\operatorname{End}_{\mathcal{O}_{\unsure{S|\mathbb{S}}}}(\unsure{\mathbb{W}|\mathbb{V}|W}(\mathcal{E}))}}\deleted{\handwritten{\unsure{=|\simeq}\mathcal{H}om_{\mathcal{O}_S}(\mathcal{E}\unsure{,}\mathcal{E})}}\handwritten{\simeq\mathcal{H}om}_{\handwritten{\mathcal{O}}_S}(\handwritten{\mathbb{V}}(\handwritten{\mathcal{E}}),\handwritten{\mathbb{V}}(\handwritten{\mathcal{E}}))$ définit une \uline{opéra}-\linebreak\uline{tion} de $G$ sur $\handwritten{\mathcal{E}}$\handwritten{,} si \inserted{\handwritten{et seulement si}} le morphisme correspondant $\handwritten{\theta}:G\handwritten{\times}_S\handwritten{\mathbb{V}}(\handwritten{\mathcal{E}})\handwritten{\longrightarrow\mathbb{V}}(\handwritten{\mathcal{E}})$ vérifie les deux conditions suivantes :

$1^\circ)$ si $\handwritten{\mu}:G\handwritten{\times}_S G\handwritten{\longrightarrow}G$ désigne le morphisme structural, le diagramme
\diagram{\begin{tikzcd}[ampersand replacement=\&]
G\handwritten{\times}_S G\handwritten{\times}_S\handwritten{\mathbb{V}}(\handwritten{\mathcal{E}})
 \arrow[r,"\operatorname{id}_G\handwritten{\times\theta}"]\arrow[d,"\handwritten{\mu\times}\operatorname{id}_{\handwritten{\mathbb{V}(\mathcal{E})}}"']
\& G\handwritten{\times}_S\handwritten{\mathbb{V}}(\handwritten{\mathcal{E}})\arrow[d,"\handwritten{\theta}"]\\
G\handwritten{\times}_S\deleted{X}\handwritten{\mathbb{V}}(\handwritten{\mathcal{E}})\arrow[r,"\handwritten{\theta}"]
\& \handwritten{\mathbb{V}}(\handwritten{\mathcal{E}})
\end{tikzcd}}
\inserted{est commutatif.}

$2^\circ)$ si $\handwritten{\varepsilon}:S\handwritten{\to}G$ est la section unité de $G$, le diagramme
\diagram{\begin{tikzcd}[ampersand replacement=\&]
S\handwritten{\times}_S\handwritten{\mathbb{V}}(\handwritten{\mathcal{E}})
 \arrow[rr,"\handwritten{\varepsilon\times}\operatorname{id}_{\handwritten{\mathbb{V}(\mathcal{E})}}"]
\& \& G\handwritten{\times}_S\handwritten{\mathbb{V}}(\handwritten{\mathcal{E}})\arrow[dl,"\handwritten{\theta}"]\\
\& \handwritten{\mathbb{V}}(\handwritten{\mathcal{E}})\arrow[ul,"\handwritten{\sim}"'] \&
\end{tikzcd}}
\inserted{est commutatif.}

Il revient\inserted{\handwritten{alors}} au même de dire que le morphisme $\handwritten{\rho}:\inserted{\handwritten{\mathcal{E}\longrightarrow}}\handwritten{\mathscr{A}}(G)\handwritten{\otimes}_{\handwritten{\mathcal{O}}_{\unsure{S|\mathbb{S}}}}\handwritten{\mathcal{E}}$ déduit de $\handwritten{\theta}$ vérifie les deux axiomes suivants :
\margin{\handwritten{$\mathcal{E}\longrightarrow\mathscr{A}(G)\otimes\mathcal{E}$}}

(CM1) Si $\handwritten{\Delta}:\handwritten{\mathscr{A}\longrightarrow\mathscr{A}\otimes\mathscr{A}}=\handwritten{\mathscr{A}}(G\handwritten{\times}_S G)$ est l'homomorphisme $\handwritten{\mathscr{A}(\theta)}$, le diagramme
\diagram{\handwritten{\begin{tikzcd}[ampersand replacement=\&]
\mathscr{A}\otimes\mathscr{A}\otimes\mathscr{A}
\& \mathscr{A}\otimes\mathcal{E}\arrow[l,"\operatorname{id}_{\mathscr{A}}\otimes\rho"']\\
\mathscr{A}\otimes\mathcal{E}\arrow[u,"\Delta\otimes\operatorname{id}_{\mathcal{E}}"]
\& \mathscr{A}\arrow[u,"\rho"']\arrow[l,"\rho"']
\end{tikzcd}}}
\inserted{est commutatif.}

(CM2) Si $\handwritten{\eta}:\handwritten{\mathscr{A}\to\mathcal{O}}_S$ est l'homomorphisme $\handwritten{\mathscr{A}(\unsure{\varepsilon|\mathcal{E}})}$, le diagramme
\diagram{\handwritten{\begin{tikzcd}[ampersand replacement=\&]
\mathcal{O}_S\otimes\mathcal{E}
\& \& \mathscr{A}\otimes\mathcal{E}\arrow[ll,"\eta\otimes\operatorname{id}_{\mathcal{E}}"']\\
\& \mathcal{E}\arrow[ul,"\sim"']\arrow[ur,"\rho"']\&
\end{tikzcd}}}
\inserted{est commutatif.}

Il revient au même de dire que le morphisme $\handwritten{\theta}_{\mathrm{af}}:G_{\mathrm{af}}\handwritten{\times}_S\handwritten{\mathbb{V}}(\handwritten{\mathcal{E}})\handwritten{\longrightarrow\mathbb{V}}(\handwritten{\mathcal{E}})$ correspondant à $\handwritten{\rho}$ définit une opération de $G_{\mathrm{af}}$ sur $\handwritten{\mathcal{E}}$ \deleted{\handwritten{\unsure{(compar I)|(Exp. I)}}} \inserted{\handwritten{(\unsure{compar|comparer} I 4.2)}}.

Par conséquent, il y a correspondance biunivoque entre opérations de $G$ sur $\handwritten{\mathcal{E}}$ et opéra-\linebreak tions de $G_{\mathrm{af}}$ sur $\handwritten{\mathcal{E}}$.

\uline{\deleted{L}\handwritten{L}emme} 11.7.- Soient $G$ un $S$-groupe quasi-compact et quasi-séparé sur $S$\handwritten{,}\inserted{\handwritten{plat et}} tel que $\handwritten{\mathscr{A}}(G)$ soit
% archive sheet number: 45
